\chapter{Attributes, properties, classes, and styles (\texttt{src/attributes.ts})}

This chapter is the ``DOM edge'' chapter.
Mounting (Chapter~5) can create elements in a few lines, but \emph{correctly configuring those elements} is where UI libraries accumulate complexity.

Relax keeps its attribute system intentionally small: a handful of rules, continuously validated by broad jsdom test coverage.

\section{Why this chapter matters}
Browsers expose multiple mechanisms that look similar but behave differently:

\begin{enumerate}
  \item \textbf{DOM attributes}: strings stored in the element's attribute map (example: \texttt{el.getAttribute('id')}).
  \item \textbf{DOM properties}: live fields on the element object (example: \texttt{(input as any).value}).
  \item \textbf{Special cased fields}: \texttt{className} and \texttt{style} have their own APIs and behavior.
\end{enumerate}

\subsection*{Diagram: the attribute write decision (attribute vs property)}
\begin{center}
\begin{tikzpicture}[
  node distance=10mm and 16mm,
  box/.style={draw, rounded corners, align=left, inner sep=6pt, fill=RelaxLight},
  decision/.style={draw, diamond, aspect=2, align=center, inner sep=2pt, fill=RelaxLight},
  arrow/.style={-Latex, thick}
]
\node[box] (input) {Input\\\texttt{name, value}};
\node[decision, right=of input] (isnull) {\texttt{value == null}?};
\node[box, above right=of isnull] (remove) {\texttt{removeAttribute(el, name)}\\(also tries property\,=\,null)};
\node[decision, below right=of isnull] (isdata) {\texttt{name.startsWith('data-')}?};
\node[box, right=of isdata] (attr) {Write attribute\\\texttt{el.setAttribute(name, String(value))}};
\node[box, below=of attr] (prop) {Write property\\\texttt{(el as any)[name] = value}};

\draw[arrow] (input) -- (isnull);
\draw[arrow] (isnull) -- node[midway, above] {yes} (remove);
\draw[arrow] (isnull) -- node[midway, left] {no} (isdata);
\draw[arrow] (isdata) -- node[midway, above] {yes} (attr);
\draw[arrow] (isdata) -- node[midway, right] {no} (prop);
\end{tikzpicture}
\end{center}

This is the core policy tested by the big matrix in \texttt{src/\_\_tests\_\_/attributes.test.ts}: most names are treated as properties, while \texttt{data-*} is always an attribute.

If you set the wrong one, you can get subtle bugs. For instance, setting \texttt{value} as an attribute on \texttt{<input>} won't behave like setting the property in many scenarios.

Relax's pragmatic rule is:

\begin{quote}
  	extbf{Rule:} if the name starts with \texttt{data-}, set an \emph{attribute}; otherwise set a \emph{property}.
\end{quote}

You can see this rule directly in \texttt{setAttribute()} in \texttt{src/attributes.ts}.

\section{Contract: \texttt{setAttributes(el, attrs)}}
The top-level API is:

\begin{quote}
	exttt{setAttributes(el: HTMLElement, attrs: ElementVNodeProps)}
\end{quote}

This function is used during mount (Chapter~5) and also during patches when props change (Chapter~8).

\subsection*{Inputs}
\begin{itemize}
  \item \texttt{el}: a real DOM element.
  \item \texttt{attrs}: Relax's element props object (roughly: \texttt{id}, \texttt{class}, \texttt{style}, and anything else).
\end{itemize}

\subsection*{Behavior}
	exttt{setAttributes()} applies attributes in three phases:

\begin{enumerate}
  \item split out \texttt{class} and \texttt{style}
  \item delete \texttt{key} (VDOM identity only)
  \item for each remaining entry: \texttt{setAttribute(el, name, value)}
\end{enumerate}

\subsection*{Diagram: \texttt{setAttributes()} pipeline (split, normalize, apply)}
\begin{center}
\begin{tikzpicture}[
  node distance=10mm and 12mm,
  box/.style={draw, rounded corners, align=left, inner sep=6pt, fill=RelaxLight},
  arrow/.style={-Latex, thick}
]
\node[box] (attrs) {Incoming props\\\texttt{\{ class, style, key, ...rest \}}};
\node[box, right=of attrs] (split) {Split\\\texttt{class/style}\\+ \texttt{rest}};
\node[box, right=of split] (sanitize) {Sanitize\\delete \texttt{key}};
\node[box, right=of sanitize] (apply) {Apply\\\texttt{setClass}\,\&\,\texttt{setStyle}\\then \texttt{setAttribute} for each rest};

\draw[arrow] (attrs) -- (split);
\draw[arrow] (split) -- (sanitize);
\draw[arrow] (sanitize) -- (apply);
\end{tikzpicture}
\end{center}

The important invariant here is that \texttt{key} never reaches the DOM, even if user code passes it.

\subsection{Important invariant: \texttt{key} is not a DOM attribute}
Relax uses \texttt{key} for diffing lists and stabilizing identity. It should never appear on the real DOM.

To enforce this, key removal happens in two places:

\begin{itemize}
  \item \texttt{src/attributes.ts}: deletes \texttt{otherAttrs.key}
  \item \texttt{src/utils/props.ts}: deletes \texttt{props.key} when extracting props/events
\end{itemize}

The duplication is intentional: it keeps each layer correct even if used in isolation.

\section{Class handling: replacement semantics}
Relax exposes \texttt{class} as either:

\begin{itemize}
  \item a string (example: \texttt{"foo bar"})
  \item an array of class tokens (example: \texttt{["foo", "bar"]})
\end{itemize}

Implementation detail (\texttt{setClass()} in \texttt{src/attributes.ts}):

\begin{itemize}
  \item it clears \texttt{el.className = ''}
  \item if \texttt{className} is a string, it assigns \texttt{el.className = className}
  \item if it's an array, it does \texttt{el.classList.add(...className)}
\end{itemize}

This makes class setting \textbf{replacement based}, not incremental diff. That's a conscious simplicity trade-off.

\subsection*{What the tests prove}
The first three tests in \texttt{src/\_\_tests\_\_/attributes.test.ts} provide the contract:

\begin{itemize}
  \item \texttt{setting a class from a string}
  \item \texttt{setting a class from an array}
  \item \texttt{updating the class} (proves replacement, not accumulation)
\end{itemize}

\section{Style handling: direct assignment}
Relax's \texttt{style} prop is a plain object:

\begin{quote}
	exttt{\{ style: \{ color: 'red', backgroundColor: 'blue' \} \}}
\end{quote}

For each \texttt{[name, value]} pair, \texttt{setAttributes()} calls \texttt{setStyle(el, name, value)}, which does a direct assignment:

\begin{quote}
	exttt{(el.style as any)[name] = value}
\end{quote}

Removal goes through \texttt{removeStyle()} which sets the field to \texttt{null}.

\subsection*{Diagram: class/style are special-cased, everything else is generic}
\begin{center}
\begin{tikzpicture}[
  node distance=10mm and 12mm,
  box/.style={draw, rounded corners, align=center, inner sep=6pt, fill=RelaxLight},
  arrow/.style={-Latex, thick}
]
\node[box] (props) {\texttt{ElementVNodeProps}};
\node[box, below left=of props] (class) {\texttt{class}\\\texttt{setClass()}};
\node[box, below=of props] (style) {\texttt{style}\\\texttt{setStyle()/removeStyle()}};
\node[box, below right=of props] (rest) {\texttt{other props}\\\texttt{setAttribute()/removeAttribute()}};
\node[box, below=of style, yshift=-2mm] (dom) {Real DOM element};

\draw[arrow] (props) -- (class);
\draw[arrow] (props) -- (style);
\draw[arrow] (props) -- (rest);
\draw[arrow] (class) -- (dom);
\draw[arrow] (style) -- (dom);
\draw[arrow] (rest) -- (dom);
\end{tikzpicture}
\end{center}

One subtle consequence: \texttt{setAttributes()} does not diff old vs new style objects. If you call it twice, the second call only assigns the fields present in the second object. it does \emph{not} remove fields that were present previously unless you explicitly remove them.

\subsection*{What the tests prove}
	exttt{setting styles} and \texttt{updating styles} assert that style fields are applied by assignment and can coexist across multiple calls.

\section{The core rule: \texttt{setAttribute(el, name, value)}}
The function \texttt{setAttribute()} is the heart of the module:

\begin{enumerate}
  \item If \texttt{value == null} \textrightarrow{} call \texttt{removeAttribute(el, name)}.
  \item Else if \texttt{name.startsWith('data-')} \textrightarrow{} call \texttt{el.setAttribute(name, String(value))}.
  \item Else \textrightarrow{} write as a property: \texttt{(el as any)[name] = value}.
\end{enumerate}

This is where Relax chooses ``property first'' semantics. The benefit is correctness for form controls and boolean-ish fields.

\subsection*{What the tests prove: the matrix}
  exttt{src/\_\_tests\_\_/attributes.test.ts} is a deliberately wide matrix of real DOM behavior.

Highlights:

\begin{itemize}
  \item \textbf{Generic element properties:}
    	exttt{hidden}, \texttt{tabIndex} (including \texttt{tabIndex} removal resulting in \texttt{-1}).
  \item \textbf{Input type=text properties:}
    	exttt{value}, \texttt{placeholder}, \texttt{disabled}, \texttt{minLength}, \texttt{size}, and more.
  \item \textbf{Checkbox / radio:}
    	exttt{checked} is a property and removal resets to \texttt{false}.
  \item \textbf{Select:}
    setting \texttt{value} to a missing option results in \texttt{''}.
  \item \textbf{Anchors/buttons/spans:}
    properties like \texttt{href}, \texttt{disabled}, and \texttt{textContent} behave like live fields.
\end{itemize}

This test suite is doing an important job: it prevents ``helpful refactors'' from changing semantics that look equivalent but aren't.

\section{Removal: \texttt{removeAttribute(el, name)} and jsdom quirks}
Removing an attribute sounds simple but isn't.

Relax tries to clear both levels:

\begin{enumerate}
  \item it attempts to set the property to \texttt{null} inside a \texttt{try/catch}
  \item it always calls \texttt{el.removeAttribute(name)}
\end{enumerate}

The \texttt{try/catch} exists because some properties throw when assigned \texttt{null} (depending on element type and environment).
When that happens, Relax logs a warning:

\begin{quote}
	exttt{Failed to set "<name>" to null on <TAG>}
\end{quote}

That warning is not a test failure; it's a visibility aid for cases where property clearing is not permitted.

\section{Props vs events: \texttt{extractPropsAndEvents} (\texttt{src/utils/props.ts})}
Relax stores event handlers under \texttt{props.on}. Before mounting or constructing a component instance, it needs to split the two.

	exttt{extractPropsAndEvents(vdom)} does exactly that:

\begin{itemize}
  \item \texttt{events}: \texttt{vdom.props.on} (defaulting to \texttt{\{\}})
  \item \texttt{props}: everything else (and it deletes \texttt{key})
\end{itemize}

This split is used in two places you've already seen:

\begin{itemize}
  \item element mounting: \texttt{addEventListeners(events, el, hostComponent)} and \texttt{setAttributes(el, props)}
  \item component mounting: \texttt{new Component(props, events, hostComponent)}
\end{itemize}

\subsection*{What the tests prove}
  exttt{src/\_\_tests\_\_/props.test.ts} is intentionally small and crisp:

\begin{itemize}
  \item empty input yields empty outputs
  \item \texttt{key} is ignored
  \item non-event fields remain in \texttt{props}
  \item \texttt{on: \{...\}} becomes \texttt{events}
\end{itemize}

Treat these as invariants: if you ever change the VNode prop shape, these tests are the first alarm bell.

\section{Design notes (trade-offs)}
Relax's attribute model is intentionally simple:

\begin{itemize}
  \item it mostly treats prop names as property names
  \item it special-cases only \texttt{data-*}, \texttt{class}, and \texttt{style}
  \item it relies on a test matrix to ensure that simplicity doesn't break real DOM behavior
\end{itemize}

If you were extending Relax, likely next steps would be:

\begin{itemize}
  \item richer SVG support (namespaces and attribute/property differences)
  \item style diff/removal semantics (tracking old style object)
  \item explicit boolean attribute reflection rules (attribute presence vs property value)
\end{itemize}

\section{Mini-exercises}
\begin{enumerate}
  \item Add a new case to the matrix in \texttt{src/\_\_tests\_\_/attributes.test.ts} for \texttt{<textarea>} value behavior.
  \item Add a test that asserts \texttt{setAttributes(el, \{ key: 'x', id: 'y' \})} does not create a \texttt{key} attribute in the DOM.
\end{enumerate}

\section{Online resources}
\begin{itemize}
  \item Attributes vs. properties (MDN, conceptual background): \href{https://developer.mozilla.org/en-US/docs/Web/API/Element}{https://developer.mozilla.org/en-US/docs/Web/API/Element}
  \item \texttt{HTMLElement.dataset} and \texttt{data-*} attributes (MDN): \href{https://developer.mozilla.org/en-US/docs/Web/API/HTMLElement/dataset}{https://developer.mozilla.org/en-US/docs/Web/API/HTMLElement/dataset}
  \item Form control properties like \texttt{value} and \texttt{checked} (MDN): \href{https://developer.mozilla.org/en-US/docs/Web/API/HTMLInputElement}{https://developer.mozilla.org/en-US/docs/Web/API/HTMLInputElement}
  \item CSSOM style API (MDN): \href{https://developer.mozilla.org/en-US/docs/Web/API/CSSStyleDeclaration}{https://developer.mozilla.org/en-US/docs/Web/API/CSSStyleDeclaration}
\end{itemize}
